The zenith angle of the sun at summer solstice will be
\[
  z_s = \phi - 23.5 = 12.5
\]
\hfill(1.5 points)

And in the winter solstice
\[
  z_w = \phi + 23.5 = 59.5
\]
\hfill(1.5 points)

Figures shows that in summer solstice we have
\[
  \tan(z_s) = \frac{x}{h} = 0.22
\]
\hfill(1.5 points)

And in the winter solstice
\[
  \tan(z_w) = \frac{D + x}{H} = 1.70
\]
\hfill(1.5 points)

Then
\[
  x = 1.70H - D = 0.60\,\mathrm{m}
\]
\hfill(2 points)
\[
  h = 2.73\,\mathrm{m}
\]
\hfill(2 points)

% FIGURE NEEDED: two room cross-sections (winter solstice and summer solstice) showing the sun ray past the Tabeshband through the window, with labels H, D, h, x, z_w, z_s, A (2009 solutions, p. 45)
